
\section{Data Analysis}

The reconstruction of the complex electric field $E(x,y)$ from the acquired digital hologram
is carried out through a fully numerical pipeline implemented in Python.
The procedure is articulated in four main stages:
spectral field recovery at the sensor plane,
spectral quality assessment,
numerical propagation along the optical axis,
and quantitative analysis of amplitude and phase.
Each stage is described in detail below.

% ─────────────────────────────────────────────────────────────────────────────
\subsection{Hologram Acquisition and FFT Computation} %per questa prima parte si possono usare le immagini del primo giorno, giusto per mostrare come funziona il codice e l'analisi dati

The hologram is recorded by a Basler acA3800-14um camera (sensor pixel pitch
$\delta_\text{sensor} = \SI{2.4}{\micro\meter}$) and saved as an OME-TIFF stack.
The raw frame is loaded as a 2D floating-point array $I(x,y)$ of size
$N_x \times N_y$ pixels.

In the off-axis configuration, the recorded intensity results from the
interference of the reference wave $R$ and the object wave $O$:
\begin{equation}
    I(x,y) \;=\; |R + O|^2
            \;=\; |R|^2 + |O|^2 + R^*O + RO^*
    \label{eq:hologram_intensity}
\end{equation}
The four terms in Eq.~\eqref{eq:hologram_intensity} produce
\textit{three distinct regions} in the Fourier domain
(Fig.~\ref{fig:fft_spectrum}):
the strong DC component (order~0, centred), the \textit{useful} sideband
(order~$+1$, which encodes the object's amplitude and phase), and its
complex-conjugate mirror (order~$-1$).
The lateral tilt between $R$ and $O$ spatially separates all three orders,
making spectral isolation of the $+1$ order possible.

The 2D Fast Fourier Transform is computed on $I(x,y)$, and the resulting
spectrum is centred with \texttt{fftshift} so that the DC term appears at
the array centre.
The logarithmic magnitude $\log(|\tilde{I}|+1)$ is displayed for
visualisation purposes.

% ─────────────────────────────────────────────────────────────────────────────
\subsection{Gated Spectrum and Field Recovery}

The function \texttt{compute\_gated\_spectrum} isolates the $+1$ order and
recovers the complex electric field $E(x,y,0)$ at the sensor plane through
four sequential steps.

\subsubsection*{Step 1 — Adaptive DC masking}
A circular mask of radius
$r_\text{exc} = 0.15\,\min(N_x,N_y)$ pixels is applied around the spectral
centre, zeroing the dominant DC component so that it does not bias the
subsequent peak search.

\subsubsection*{Step 2 — Sub-pixel peak localisation}
The global maximum in the masked spectrum is located, and its position is
refined to sub-pixel accuracy by computing the \textit{centre of mass} of a
$5{\times}5$~pixel neighbourhood using
\texttt{scipy.ndimage.center\_of\_mass}.
This yields the peak coordinates $(i_\text{peak},\,j_\text{peak})$ with
sub-pixel precision.

\subsubsection*{Step 3 — Rectangular spectral cropping}
A window of size $\Delta X \times \Delta Y$ pixels is extracted from the
centred spectrum, centred on the $+1$ peak.
The vertical extent is rescaled proportionally to the image aspect ratio:
\begin{equation}
    \Delta Y \;=\; \Delta X \cdot \frac{N_x}{N_y}
\end{equation}
The choice of $\Delta X$ sets a trade-off between spatial resolution
and noise rejection (see Section~\ref{sec:resolution_budget}).

\subsubsection*{Step 4 — Complex field recovery via IFFT}
The inverse FFT of the cropped window (after inverse-shifting with
\texttt{ifftshift}) yields the complex electric field at the sensor plane:
\begin{equation}
    E(x,y,0) \;=\; \mathcal{F}^{-1}\!\left\{\,\tilde{I}_{+1}(f_x,f_y)\,\right\}
    \label{eq:field_ifft}
\end{equation}

\subsubsection*{Physical parameters}
The object-plane pixel pitch is determined by the optical magnification
$M$ of the microscope objective:
\begin{equation}
    \delta_\text{obj} \;=\; \frac{\delta_\text{sensor}}{M}
\end{equation}
The parameters used throughout the analysis are summarised in
Table~\ref{tab:setup_params}.

\begin{table}[h!]
    \centering
    \caption{Physical setup parameters.}
    \label{tab:setup_params}
    \begin{tabular}{lcc}
        \hline
        \textbf{Parameter} & \textbf{Symbol} & \textbf{Value} \\
        \hline
        Sensor pixel pitch     & $\delta_\text{sensor}$ & \SI{2.4}{\micro\meter} \\
        Optical magnification  & $M$                    & $12\times$ \\
        Object-plane pixel pitch & $\delta_\text{obj}$  & \SI{0.2}{\micro\meter} \\
        Laser wavelength       & $\lambda$              & \SI{635}{\nano\meter}\ (HeNe) \\
        Numerical aperture     & NA                     & $0.3$ \\
        Gating window          & $\Delta X$             & $300\;\text{px}$ \\
        \hline
    \end{tabular}
\end{table}

% ─────────────────────────────────────────────────────────────────────────────
\subsection{Resolution Budget: Digital vs.\ Optical Limit}
\label{sec:resolution_budget}

The size of the gating window $\Delta X$ directly controls the spatial
resolution of the reconstructed field.
The spatial frequency step and the resulting digital resolution along each
axis are:
\begin{equation}
    \delta f_x \;=\; \frac{1}{N_x\,\delta_\text{obj}}, \qquad
    B_x \;=\; \Delta Y \cdot \delta f_x, \qquad
    \delta x_\text{dig} \;=\; \frac{1}{B_x}
    \label{eq:digital_resolution}
\end{equation}
and analogously along $y$.

The theoretical optical resolution is given by the Rayleigh criterion for
coherent illumination:
\begin{equation}
    \delta_\text{optical} \;=\; \frac{0.61\,\lambda}{\mathrm{NA}}
    \label{eq:rayleigh}
\end{equation}
An equivalent \textit{experimental} numerical aperture can be derived by
inverting Eq.~\eqref{eq:rayleigh}:
\begin{equation}
    \mathrm{NA}_\text{exp} \;=\; \frac{0.61\,\lambda}{\delta_\text{dig}}
\end{equation}
If $\delta_\text{dig} > \delta_\text{optical}$, the gating window is too
narrow and is discarding spatial frequencies physically present in the signal.
The minimum window size required to capture the full optical bandwidth is:
\begin{equation}
    \Delta X_\text{min} \;=\; \frac{2\,f_\text{max}}{\delta f_y}, \qquad
    f_\text{max} \;=\; \frac{\mathrm{NA}}{0.61\,\lambda}
\end{equation}
This diagnostic check is systematically performed by the notebook and
issues a warning whenever the chosen $\Delta X$ is insufficient.

% ─────────────────────────────────────────────────────────────────────────────
\subsection{Spectral Masking}

Two complementary masks are applied to the gated spectrum in order to
suppress noise and residual artefacts.

\paragraph{Elliptic outer mask.}
Only spectral components within an ellipse of semi-axes
$a = \Delta Y/2$ and $b = \Delta X/2$ (consistent with the physical extent of
the cropping window) are retained; all components outside are set to zero.

\paragraph{Central-peak removal.}
A disk of fixed radius $r_0 = 13$~px, centred on the dominant spectral peak
inside the ellipse, is zeroed.
This eliminates the residual DC contribution or any strong undiffracted
component that would otherwise corrupt the reconstructed phase.

The combined binary mask is:
\begin{equation}
    M(f_x,f_y)
    \;=\;
    M_\text{ellipse}(f_x,f_y)
    \;\wedge\;
    \overline{M}_\text{circle}(f_x,f_y)
    \label{eq:combined_mask}
\end{equation}

Both the real-valued (modulus) and the complex-valued gated spectra are
processed with this mask.
The masked complex spectrum is retained for optional alternative
reconstructions; the main propagation pipeline uses the \textit{unmasked}
field $E(x,y,0)$.

% ─────────────────────────────────────────────────────────────────────────────
\subsection{Frequency-Weighted Sharpness Metric}

A spectral sharpness metric $S_\rho$ is introduced to quantify the
high-frequency content of the reconstructed field — and thereby to assess
the degree of focus.
For each active pixel $(f_x,f_y)$ in the masked gated spectrum, the
normalised radial distance from the ellipse centre is:
\begin{equation}
    \rho(f_x,f_y)
    \;=\; \sqrt{\left(\frac{f_x - f_{x,0}}{a}\right)^2
               +\left(\frac{f_y - f_{y,0}}{b}\right)^2}
\end{equation}
The weighted sharpness is then defined as:
\begin{equation}
    S_\rho
    \;=\;
    \frac{\displaystyle\sum_{\text{valid}}
          \log\!\left(1+|\tilde{E}|\right)\cdot\rho}
         {\displaystyle\sum_{\text{valid}}\rho}
    \label{eq:sharpness}
\end{equation}
A high value of $S_\rho$ indicates that spectral power is concentrated
towards the periphery of the ellipse, corresponding to a sharp, well-focused
reconstruction.
A low value indicates that energy is confined to low spatial frequencies,
signalling defocus.
This metric can be evaluated as a function of propagation distance $z$
(Section~\ref{sec:propagation}) to automate the search for the best-focus
plane.

% ─────────────────────────────────────────────────────────────────────────────
\subsection{Numerical Propagation — Angular Spectrum Method}
\label{sec:propagation}

\subsubsection*{Theoretical framework}
Given the complex field $E(x,y,0)$ at the sensor plane, the field at an
arbitrary axial distance $z$ is obtained by multiplying in the Fourier domain
by the \textit{Angular Spectrum} transfer function:
\begin{equation}
    E(x,y,z)
    \;=\;
    \mathcal{F}^{-1}\!\left\{
        H(f_x,f_y;\,z)\;\cdot\;\tilde{E}(f_x,f_y)
    \right\}
    \label{eq:asm}
\end{equation}
where
\begin{equation}
    H(f_x,f_y;\,z) \;=\; e^{-i k_z z},
    \qquad
    k_z \;=\; \sqrt{k_0^2 - 4\pi^2(f_x^2+f_y^2)},
    \qquad
    k_0 \;=\; \frac{2\pi}{\lambda}
    \label{eq:propagator}
\end{equation}
Spatial frequencies satisfying $f_x^2+f_y^2 > (1/\lambda)^2$ correspond to
\textit{evanescent waves}: $k_z$ becomes purely imaginary and their
contribution decays exponentially away from the source plane.
These components are suppressed by a circular bandpass mask of cutoff radius
$f_0 = 1/\lambda$.
Numerical stability for evanescent frequencies is ensured by casting
$k_z^2$ to \texttt{complex} before taking the square root, which prevents
\texttt{NaN} propagation.

\subsubsection*{Implementation}
The propagation routine is factored into three modular functions:

\begin{itemize}
    \item \texttt{circle\_mask($\mathbf{FF}_x$,$\mathbf{FF}_y$,$f_0$)} —
          returns a binary mask equal to~1 for propagating frequencies
          ($f_x^2+f_y^2 \leq f_0^2$) and~0 elsewhere;
    \item \texttt{propagator($k_0$,$\mathbf{FF}_x$,$\mathbf{FF}_y$,$f_0$,$z$)} —
          computes the transfer function $H$ of Eq.~\eqref{eq:propagator};
    \item \texttt{step\_propagation($E_0$,\,freq,\,$f_0$,$k_0$,$z$)} —
          executes the full pipeline:
          FFT$_2 \to$ multiply by $H \to$ IFFT$_2$.
\end{itemize}

\subsubsection*{Numerical refocusing}
Starting from $E(x,y,0)$, the field is propagated over a discrete set of
axial planes:
\begin{equation}
    z \;\in\; [-24,\;+24]\;\si{\micro\meter},
    \quad \Delta z = \SI{2}{\micro\meter}
    \quad (25\;\text{planes total})
\end{equation}
For each plane, $|E(x,y,z)|$ is computed and displayed.
The best-focus depth $z^*$ is identified as the plane that maximises
the sharpness metric $S_\rho$ defined in Eq.~\eqref{eq:sharpness}.

% ─────────────────────────────────────────────────────────────────────────────
\subsection{Amplitude and Phase Reconstruction}

From the recovered complex field $E(x,y,0) = |E|\,e^{i\phi}$,
three quantitative maps are extracted.

\paragraph{Amplitude.}
\begin{equation}
    A(x,y) \;=\; |E(x,y,0)|
\end{equation}
Bright regions correspond to areas of high scattering or absorption.

\paragraph{Wrapped phase.}
\begin{equation}
    \phi_w(x,y) \;=\; \arg\!\left(E(x,y,0)\right) \;\in\; [-\pi,\,+\pi]
\end{equation}
Computed via \texttt{numpy.angle}.
Wherever the accumulated phase exceeds $2\pi$, the $\arctan$ function
introduces discontinuities (phase wrapping) visible as abrupt $2\pi$~jumps.

\paragraph{Unwrapped phase.}
\begin{equation}
    \phi(x,y) \;=\; \mathcal{U}\!\left[\phi_w(x,y)\right]
\end{equation}
Computed with \texttt{skimage.restoration.unwrap\_phase}, which implements
a 2D path-following phase-unwrapping algorithm.
The result is a continuous, quantitative phase map related to the optical
path length through the sample:
\begin{equation}
    \phi(x,y)
    \;=\; \frac{2\pi}{\lambda}\,\Delta n(x,y)\,t(x,y)
    \label{eq:opd}
\end{equation}
where $\Delta n$ is the refractive-index contrast with respect to the
surrounding medium and $t$ is the local sample thickness.
Equation~\eqref{eq:opd} is the primary relation linking the measured phase
to the physical properties of the specimen, enabling quantitative
morphological and optical characterisation.

% ─────────────────────────────────────────────────────────────────────────────
\subsection{Processing Pipeline Summary}

The complete data-analysis workflow is schematised below.

\begin{figure}[h!]
\centering
\begin{tikzpicture}[
    node distance   = 0.55cm and 1.2cm,
    box/.style      = {draw, rounded corners=3pt, fill=blue!8,
                       minimum width=5.2cm, minimum height=0.65cm,
                       font=\small, align=center},
    side/.style     = {draw, rounded corners=3pt, fill=orange!12,
                       minimum width=4.0cm, minimum height=0.65cm,
                       font=\small, align=center},
    arr/.style      = {->, thick},
    sarr/.style     = {->, dashed, gray}
]

\node[box] (load)   {Raw hologram $I(x,y)$ — TIFF/OME-TIFF};
\node[box, below=of load]  (fft)    {2D FFT $\to$ centred spectrum $\tilde{I}(f_x,f_y)$};
\node[box, below=of fft]   (gate)   {Adaptive DC mask + sub-pixel peak search};
\node[box, below=of gate]  (crop)   {Rectangular crop $\Delta X \times \Delta Y$ on order $+1$};
\node[box, below=of crop]  (ifft)   {IFFT $\to$ complex field $E(x,y,0)$};

\node[side, right=of crop]  (rescheck) {Resolution budget\\(digital vs.\ Rayleigh)};
\node[side, right=of ifft]  (mask)     {Elliptic mask +\\central-peak removal};
\node[side, below=of mask]  (sharp)    {Weighted sharpness\\$S_\rho$};

\node[box, below=of ifft]  (asm)    {Angular Spectrum propagation\\$E(x,y,z),\; z\in[-24,+24]\,\mu\text{m}$};
\node[box, below=of asm]   (maps)   {Amplitude $|E|$\quad Wrapped $\phi_w$\quad Unwrapped $\phi$};

\draw[arr] (load)  -- (fft);
\draw[arr] (fft)   -- (gate);
\draw[arr] (gate)  -- (crop);
\draw[arr] (crop)  -- (ifft);
\draw[arr] (ifft)  -- (asm);
\draw[arr] (asm)   -- (maps);

\draw[sarr] (crop)  -- (rescheck);
\draw[sarr] (ifft)  -- (mask);
\draw[sarr] (mask)  -- (sharp);
\draw[sarr] (sharp) -- (asm);

\end{tikzpicture}
\caption{Schematic of the digital holographic reconstruction pipeline.
         Solid arrows indicate the main data flow; dashed arrows indicate
         diagnostic/optional branches.}
\label{fig:pipeline}
\end{figure}


